% ===================================================================
%  THE VOLUME BEGINS AND ENDS AT THE PRINTED COVER, green paper
%  (leaves 3 and 4 at the head, 237 and 238 at the tail). The BINDING
%  is not transcribed: leaves 1 and 240 are the two boards of a
%  publisher's case in brown buckram with a gilt title, leaves 2 and
%  239 the endpapers that line them. They were surveyed -- the front
%  board carries three gilt lines and a thin rule, the back is bare --
%  but they belong to the copy, not to the edition, and the
%  transcription stops at what the printer set.
%  The cover's paper is known by measurement: median RGB 177/163/139
%  at leaves 3 and 4, 177/169/152 and 176/175/161 at 237 and 238,
%  against 237/210/184 at the paper of the body of the book.
% ===================================================================

% ===================================================================
%  LEAF 3 -- THE PRINTED COVER (front cover).
%  The volume therefore carries both a publisher's binding (leaves 1
%  and 240) AND the printed cover, kept at the binding.
%
%  CALIBRATION: the TOP EDGE OF THE PAPER, taken at the middle of the
%  scan's shadow ramp -- y = 38 px, +- 4 px -- then the constant
%  C = 5.84 mm of leaf 7. Total uncertainty +- 1.6 mm on the position
%  of the whole; the internal distances are measured to the pixel.
%  Deskewed by -1.60 degrees.
%
%  THE LEFT EDGE IS CROPPED BY THE SCANNING: at x=0 the luminance is
%  213 and FALLS inwards -- there is no shadow ramp, so the paper
%  continues out of frame. Some 33 px = 2.8 mm are missing on the left,
%  and the measurable width is only 108 mm against the volume's 116.
%  THE AXIS OF CENTRING IS THEREFORE ESTIMATED, NOT MEASURED.
%
%  « LINGUO INTERNACIONA » OVERRUNS THE MEASURE: 1148 px against the
%  volume's 1083, that is 5.6 mm too far. It is the same overrun as on
%  the title page, and \VUcentreA lets it overrun instead of folding
%  it, the line being laid in an \hbox and not in a paragraph.
%
%  THE VIGNETTE is the emblem of Ido -- a six-pointed star carrying
%  IDO, circled by the legend LINGUO INTERNACIONA DI LA DELEGITARO.
%  Like the portrait of leaf 7 and the fleuron of folio 232, it is CUT
%  FROM THE SCAN: rot[1207:1456, 482:723], 241 x 249 px =
%  20.40 x 21.08 mm.
%
%  THE LAST TWO LINES ARE ONLY ONE: « Imprimita en Luxemburg. » and
%  « Preco : 2 Suisa Franki. » share the baseline y=1893, one flush
%  left, the other flush right. Hence \VUcaseA, which lays the two on
%  a given measure.
%
%  TWO WEIGHTS NOT DECIDED, and said to be so:
%   -- « DI LA »: stem/capital ratio 0.250, as high as the bold IDO,
%      but on a capital of 20 px where one pixel of stem shifts the
%      ratio by 0.05. Enlarged x6, it is an ordinary roman a little
%      fattened by the spread of the ink. Set roman, not proved.
%   -- « 1925 »: 0.174, between the roman (0.154) and the semibold
%      (0.192) of the same page. Set roman, without certainty.
% ===================================================================
\begin{VUpage}[3]{}
\VUcentreA{16.25mm}{10.63pt}{109}{\VUgras{L. de BEAUFRONT}}
\VUfiletA[1.92pt]{23.79mm}{12.70mm}
\VUcentreA{37.51mm}{19.13pt}{112}{\VUetroit{0.668}{KOMPLETA GRAMATIKO DETALOZA}}
\VUcentreA{50.29mm}{7.09pt}{276}{DI LA}
\VUcentreA{61.04mm}{23.74pt}{29}{LINGUO INTERNACIONA}
\VUcentreA{76.54mm}{41.09pt}{4}{\VUgras{IDO}}
\VUimageA{104.72mm}{22.38mm}{ornamenti/vinieto-3}
\VUcentreA{147.32mm}{8.86pt}{0}{\VUgras{Editerio : MEIER-HEUCKE}}
\VUcentreA{152.23mm}{9.21pt}{0}{\VUgras{Esch-Alzette} \VUgras{\textit{(Luxemburgia)}}}
\VUfiletA[0.96pt]{156.97mm}{2.62mm}
\VUcentreA{160.10mm}{8.15pt}{80}{1925}
\VUcaseA{164.84mm}{102.53mm}{9.21pt}{Imprimita en Luxemburg.}{Preco : 2 Suisa Franki.}
\end{VUpage}

% ===================================================================
%  LEAF 4 -- THE INSIDE FRONT COVER: « QUO ESAS IDO? »
%  It is not a blank verso, as the journal said: it is a FULL page,
%  45 lines, measure 1080 px = 91.44 mm, that is the volume's to
%  within 0.25 mm.
%
%  THE LINE PITCH IS 33 px (2.79 mm), NOT 41: it is the leading of the
%  volume's NOTES, and exactly that of leaves 237-238, the publisher's
%  announcements, already set with \VUpageSerre{8.26pt}{7.95pt}. The
%  capital height of the body is 23-24 px. The paragraph indent is
%  34 px = 2.88 mm, not the 4.6 mm of \VUalinea: it is laid in the
%  page.
%
%  TWO DISTINCT PARAGRAPH WHITES, and this is the most useful
%  measurement on the page: in the upper section, an excess of 7-8 px
%  = 1.69 to 1.93 pt, that is \VUblancAlinea (1.86 pt); in the
%  « Opinioni » section, an excess of 12 px = 2.89 pt.
%
%  THE CONVENTION OF ENRICHMENT IS REVERSED, and kept as it stands:
%  section 3.3 of the journal gives bold = an Ido form as a headword,
%  italic = a word of another language. Here the IDO words used as
%  headwords -- Komitato, Delegitaro, Progreso, un, Ido -- are set in
%  ITALIC. That is consistent: this page is not part of the volume's
%  body, it is a publisher's text.
%
%  « 1,250 » with the thousands comma is indeed what is printed.
%
%  THE PAGE STOPS IN THE MIDDLE OF A SENTENCE -- « ...ta principi
%  aplikesis kun rigoro » -- and the sequel is nowhere in the scan:
%  the inside back cover was not bound in. The last line therefore
%  carries \nl and not \parplein: it is not the end of a paragraph.
% ===================================================================
\begin{VUpage}[4]{}
\VUcentreA{16.04mm}{14.17pt}{45}{\VUgras{QUO ESAS IDO?}}
\VUfiletA[1.44pt]{23.66mm}{8.13mm}
\VUcentreA{118.57mm}{9.92pt}{73}{\VUgras{Opinioni pri Ido.}}
\vspace*{4.61mm}
\VUpageSerre{8.26pt}{7.95pt}
\setlength{\parindent}{2.88mm}

Ido esas l'idiomo internaciona, helpanta e neutra, qua rezultas\nl
de la selekto e labori dil « \textit{Délégation} » (delegitaro).

\VUblanc{1.86pt}\textit{La Délégation pour l'adoption d'une langue auxiliaire interna}\cc
\textit{tionale}, fondita en 1901 (januaro), esis ensemblo de personi delegita\nl
da 310 kongresi o societi omnalanda ed omnaskopa, por examenar la\nl
multa projeti o sistemi di linguo internaciona, e por adoptar la\nl
maxim bona. Til 1907 ol recevabis 1,250 aprobanta signaturi de\nl
membri di Akademii od Universitati, qui tale montris sua intereso\nl
al entreprezo.

\VUblanc{1.86pt}Lore ol elektis (junio 1907) internaciona \textit{Komitato} di ciencisti e\nl
linguisti partikulare kompetenta, qui, pos exameno di omna projeti\nl
ancien o recenta di Linguo internaciona ante e dum 18 kunsidi,\nl
adoptis fine Esperanto kun diversa modifiki. Ta modifiki studiita,\nl
kompletigita per laboro komuna, publik ed internaciona, de 1908 til\nl
1914, en la revuo \textit{Progreso}, havis kom rezultanto la nuna « Ido ».\nl
Ica konseque ne esas, quale Esperanto, la simpla produkturo di \textit{un}\nl
homo.

\VUblanc{1.86pt}Esas notenda, ke la Komitato dil \textit{Delegitaro} kontenis Esperantisti\nl
e mem la Prezidero ipsa di la « Lingva Komitato » esperantista.\nl
Or la decido adoptanta Esperanto « \textit{kun la rezervo di ula modifiki} »,\nl
esis unanima. Se, pose, skismo eventis, la responso ne esas imputebla\nl
a l'Idistaro, ma ad ilti qui obliviis promisi, voti o signaturo.

\VUblanc{1.86pt}La Komitato dil Delegitaro (la maxim kompetenta e maxim sen\cc
partisa qua ultempe kunvenis pri linguo internaciona) tale judi\cc
kesas da S\textsuperscript{ro} \textsc{Gaston Moch}, qua supleis en lu rektoro \textsc{Emile Boirac},\nl
tre eminenta chefo Esperantista : « La lingui Germana e Franca\nl
havis ibe tri reprezentanti, la Angla du, e la lingui Dana, Hispana,\nl
Greka e Hungariana, un; la maxim diversa konocaji reprezentesis,\nl
ed on remarkis en lu nome quar reputata filologi, S\textsuperscript{ri} \textsc{Baudouin}\nl
\textsc{de Courtenay}, \textsc{Jespersen}, \textsc{Lambros} e \textsc{Schuchardt}. »

\VUblanc{18.32pt}« \textit{La gramatiko di Ido satisfacas plu bone la postuli di linguo}\nl
\textit{internaciona kam la gramatiko di Esperanto...} » (A. \textsc{Meillet}, pro\cc
fesoro pri linguistiko en Collège de France, \textit{Revue Critique}, 11 di\nl
marto 1911.)

\VUblanc{2.89pt}« Esas eroro nepardonebla institucar, quale agis Esperanto,\nl
distingo pri l'akuzativo e la nominativo, distingo qua jenos omna\nl
individui di linguo romanal od Angla, e qua esas neutila a la ceteri. »\nl
(La sam autoro, \textit{Revue Critique}, 11 di marto 1911.)

\VUblanc{2.89pt}« Cetere, esas facila procedar plu logikoze e, konseque plu satis\cc
facante e klare, en la vortifado kam agas lo Esperanto.

\VUblanc{2.89pt}« Icon montris la kreinti di \textit{Ido}, linguo fondita sur la sama\nl
principi kam Esperanto, ma ube ta principi aplikesis kun rigoro\nl
\end{VUpage}

% ===================================================================
%  Leaves 5 and 6 -- preliminaries (folios 1 and 2)
%
%  Pages of apparatus: each line is placed at its MEASURED ORDINATE
%  and takes up no height; no line displaces another.
%  Measurements by tools/apparatus.py (leaf 5):
%    line        y (px)  y (mm)  cap (mm)  width (mm)   size    l.sp.
%    KOMPLETA...    754   63,84    2,540      60,20      10,63pt      2
%    DI LA          919   77,79    1,439       7,03       6,02pt    224
%    LINGUO...     1063   89,98    3,852      89,15      16,12pt    117
% ===================================================================

% --- Leaf 5 (folio 1): half title ----------------------------------
\begin{VUpage}[5]{}
\VUcentreA{63.84mm}{10.63pt}{-5}{KOMPLETA GRAMATIKO DETALOZA}
\VUcentreA{77.79mm}{6.02pt}{264}{DI LA}
\VUcentreA{89.98mm}{16.12pt}{102}{LINGUO INTERNACIONA IDO}
\end{VUpage}

% --- Leaf 6 (folio 2): blank ----------------------------------------
\begin{VUpage}[6]{}
\end{VUpage}

% -------------------------------------------------------------------
%  LEAF 7 (folio 3) -- THE TITLE PAGE, and the plate.
%  Ten elements, all of them apparatus: not one line of running text.
%  The surveyor found them centred on one optical axis to within two
%  pixels, and showed that the detector had lost two of them (LINGUO
%  INTERNACIONA IDO, whose 77 px capitals overrun the height filter,
%  and DA, 41 px wide) while inventing six, which are artefacts of the
%  halftone screen.
%
%  THE PAGE HAS NEITHER FOLIO NOR LINE OF BODY: the ordinary rule --
%  lay the ordinate relative to the first line of body -- has nothing
%  to hold on to. The first state of this page was therefore laid in
%  FLOW, with a head white chosen by eye. It was wrong: comparing the
%  render with the facsimile showed that EVERY white in it was too
%  short, the page compressing by 139 px from the top downwards.
%
%  It is taken up again here in ABSOLUTE ordinates, measured on the
%  scan deskewed by 0.20 degrees:
%    ordinate = (top_of_capitals - edge_of_paper) x 25,4/300 + C
%  where C = 5.84 mm is the part of the head margin the scanning crops
%  away. C is not a free figure: it is measured over 33 pages of text
%  where the first line of body falls, by construction, at
%  \VUmargeSup; C = 24,30 - (y_line - y_paper) x 25,4/300. The median
%  is 5.84 mm, the standard deviation 1.5 mm apart from three title
%  pages. IT IS THE ONLY APPROXIMATE QUANTITY ON THIS PAGE, and it
%  bears on the position of the block as a whole, not on its internal
%  distances, which are measured exactly.
%
%  The top of the facsimile's capitals is corrected for the spread of
%  the ink: (height_fac - height_set)/2, that is 0.5 to 6 px.
%
%  THE LINE « LINGUO INTERNACIONA IDO » OVERRUNS THE MEASURE: 1164 px
%  against the volume's 1083. It is measured three times; it is not
%  the fore-edge of the neighbouring leaf. \VUcentreA lays it in an
%  \hbox and not in a paragraph: it overruns as in the facsimile
%  instead of folding.
% -------------------------------------------------------------------
\begin{VUpage}[7]{}
\VUcentreA{10.71mm}{18.68pt}{0}{\VUetroit{0.702}{KOMPLETA GRAMATIKO DETALOZA}}
\VUcentreA{20.36mm}{6.30pt}{362}{DI LA}
\VUcentreA{27.94mm}{24.10pt}{-79}{LINGUO INTERNACIONA IDO}
\VUcentreA{40.00mm}{5.90pt}{242}{DA}
\VUcentreA{45.17mm}{8.80pt}{282}{\VUgras{L. de BEAUFRONT}}
\VUcentreA{49.06mm}{5.90pt}{222}{PREZIDERO DI LA FRANC-IDISTA SOCIETO}
\VUplancheA{55.88mm}
\VUcentreA{148.76mm}{5.90pt}{124}{Markezo L. \textsc{de Beaufront}, precipua autoro di \textit{Ido}}
\VUcentreA{157.31mm}{7.60pt}{148}{Imprimerie SOLIMPA, Luxembourg et Esch/Alzette}
\VUcentreA{160.99mm}{7.60pt}{-21}{1925}
\end{VUpage}

% ===================================================================
%  LEAF 8 (folio 4) -- THE VERSO OF THE TITLE PAGE: the KONSTATO of
%  the Akademio Idista, and its signature.
%
%  The 266,473 px of ink announced at threshold 140 are misleading:
%  only 54,934 are typographic ink, the rest is the scanner's black
%  ground around the paper (paper: y 37-2109, x 34-1310).
%
%  THE PAGE IS DELIBERATELY OPEN: the pitch of the baselines is
%  57.667 px = 13.89 pt, against 9.99 pt in the running body, at an
%  UNCHANGED SIZE (x-height 20-21 px, that of leaf 10). Hence
%  \VUinterlignePage.
%
%  CALIBRATION: the page has neither a folio nor a line of body at
%  \VUmargeSup -- its body begins at 54.27 mm because a title and a
%  deep head margin come first. It is therefore leaf 7's calibration
%  that serves: (y - edge_of_paper) x 25,4/300 + C, edge at y=37,
%  C=5.84 mm.
%  THE SURVEYOR'S RESERVATION: his own local sample (leaves 10, 12 and
%  16, the only neighbours where the detector yields both folio AND
%  first line) gives C = 6.27 mm. The whole block may therefore be
%  0.43 mm too high. Three samples against thirty-three: the project's
%  constant is kept, and the discrepancy is stated.
%
%  THE SIGNATURE BLOCK IS NOT CENTRED ON THE MEASURE: its four lines
%  have their middles at 900.5 / 901.5 / 900.5 / 901.0 px for a middle
%  of the measure at 706.5 -- an axis of its own, shifted 16.46 mm to
%  the right, constant to within one pixel. Hence \VUcentreDA.
%  The date is set flush LEFT, 80 px = 6.77 mm from the margin.
%
%  WEIGHT: the title is semibold (the stem of the K 5.94 px against
%  4.78 at the roman K of « Komitato »), but its capital height,
%  29 px, is that of the RUNNING BODY and not that of a section title
%  (44 px at leaf 9). In the body, only the quoted title « Kompleta
%  Gramatiko detaloza di la linguo internaciona Ido » is bold: stems
%  4.39 to 4.62 px against 3.64 to 3.95 at the roman of the same line.
%  The word « naciona » was the only doubtful case -- the density of
%  bold, the stem of a roman: enlarged x6 it is of the same weight as
%  « quala » which follows it, its density being inflated by the
%  absence of an ascender and a descender in the band of measurement.
%  It is therefore ROMAN.
%
%  A PARTICULARITY KEPT: the same phrase appears twice running with
%  two opposite treatments -- in bold and lower case as the TITLE OF A
%  WORK, then in roman with an initial capital as the NAME OF A
%  LANGUAGE. That is not a slip, it is the distinction intended.
%  The paragraph indent of this page is 42 px = 3.56 mm, and not the
%  4.6 mm of \VUalinea; it is not forced.
% ===================================================================
\begin{VUpage}[8]{}
\VUinterlignePage{13.89pt}
\VUcentreA{44.53mm}{10.06pt}{119}{\VUgras{\VUetroit{0.793}{KONSTATO.}}}
\VUcentreDA{91.52mm}{16.46mm}{10.2pt}{-13}{\textit{La prezidanto :}}
\VUcentreDA{96.43mm}{16.46mm}{10.2pt}{-66}{F. \textsc{Schneeberger.}}
\VUcentreDA{102.70mm}{16.46mm}{10.2pt}{-13}{\textit{La vice-sekretario :}}
\VUcentreDA{107.61mm}{16.46mm}{10.2pt}{-86}{J. \textsc{Guignon.}}
{\renewcommand{\VUsignatureX}{6.77mm}%
 \VUsignature{117.43mm}{\fontsize{10.2pt}{10.2pt}\selectfont\textls[-16]{\textit{15 julio 1925.}}}}
\vspace*{29.97mm}

L'Akademio Idista konstatas, ke la linguala formo pri\cc
zentita en ica \VUgras{Kompleta Gramatiko detaloza di la linguo}\nl
\VUgras{internaciona Ido} esas la oficala formo di la Linguo inter\cc
naciona \VUgras{Ido}, quala ol rezultis de la labori di la Linguala\nl
Komitato di la Delegitaro por l'adopto di Linguo Interna\cc
ciona, di lua konstanta komisitaro e la decidi dil Akademio\nl
Idista.
\end{VUpage}

% -------------------------------------------------------------------
%  LEAVES 9 to 12 (folios 5 to 8) -- THE AVERTO.
%  Four pages surveyed and checked separately, one per surveyor.
%  No bold in the body nor in the notes of the four pages: the italic
%  is everywhere THINNER than the roman beside it, measured stem by
%  stem against a roman word on the SAME line.
% -------------------------------------------------------------------

% --- Leaf 9 (folio 5): the title AVERTO ----------------------------
%  NO PRINTED FOLIO on this page: looked for at threshold 165 across
%  the whole width, at \VUfolioY as at the foot. The argument stays
%  empty.
%  NO RULE under the title: looked for between y=471 and y=557 at
%  thresholds 150, 170, 185 and 200 -- nothing but the grain of the
%  paper.
%  The title « Averto » is in semibold LOWER CASE with an initial
%  capital, capital height 44 px (3.73 mm): it is the only title in
%  the volume that is not in capitals. Stem 8 px for an x-height of
%  27 px (0.296) against 4 px for 21 px at the page's roman (0.190).
%  Its reference « (1) » is THIN, at the size of the text, raised by
%  13 px.
%  Note rule: y=1613, 213 px, set on the margin (x0=90 for a first
%  line of note at x0=89) -- 123.36 mm.
%  MIND: the ordinate is computed on the TITLE, not on a line of body.
%  Recomputed on the first line of body it would give 113.54 mm, that
%  is 10 mm too high.
%  The block of note (2) is laid 23 px (2 mm) further right than that
%  of note (1) and than the rule -- checked on the RAW image, it is
%  not an artefact of the deskewing. Reported, not reproduced.
% -------------------------------------------------------------------
\begin{VUpage}[9]{}
\VUnotes{122.78mm}{%
(1) L'unesma parto di ca averto esas nur la vortopa tradukuro\nl
di olta qua aparis France en 1908, lor l'edito di \textit{Grammaire complète}.

(2) \textit{La Délégation pour l'adoption d'une Langue auxiliaire inter}\cc
\textit{nationale}, fondita en 1901 (januaro), recevabis til 1907 la adhero\nl
di 310 societi omnalanda, e la aprobo di 1250 membri di Akademii\nl
ed Universitati. Ol elektis internaciona Komitato di ciencisti e lin\cc
guisti, qua pos exameno di omna projeti anciena e nova di mondo\cc
linguo, facis la decido quan on raportas hike.%
}

\VUtitre{15.66pt}{-8}{\VUetroit{0.822}{Averto}\VUtitreAppel[1.83mm]{(1)}}
\VUsaut{5.0mm}

Ica verko esas ri-imprimuro grande modifikita di la\nl
\textit{Grammaire complète} qua, imprimita nur ye 200 exempleri\nl
e ne publikigita, livresis al dispono dil Komitato di la Dele\cc
gitaro en la 15-ma dio di oktobro 1907, kun \textit{Exercaro} e\nl
\textit{Vortaro} specimena, sub la pseudonimo « Ido » (2).

On savas, ke, pos 18 kunsidi, ek qui 5 konsakresis al dis\cc
kuto komparanta di ta laboruro e di Esperanto, unanime\nl
la Komitato decidis « adoptar Esperanto principe... kun la\nl
rezervo di ula modifiki exekutenda dal permananta Komi\cc
sitaro, segun la sinso determinita dal konkluzi di la Raporto\nl
dil sekretarii e dal projeto da Ido ». Ja la Komisitaro indi\cc
kabis ula modifiki facenda en ta projeto. Sualatere la Komi\cc
sitaro permananta submisis lu a detaloza revizo, egardante\nl
plura emendi propozita; e la rezultajo di ta laboro men\cc
cionesis en Raporto publikigita en l'unesma numero di la\nl
revuo \textit{Progreso} (marto 1908).

La \textit{Grammaire complète} quan ni publikigas prezente esas\nl
strikte konforma al decidi di la Komitato e dil Komisitaro\nl
permananta; do ol ne plus esas, same kam la linguo ipsa,\nl
verko pure individuala; ol esas la final ed oficala rezultajo\nl
di la deliberi dil Komitato elektita reguloze dal \textit{Delegitaro},\nl
pos sep yari de propagado por l'ideo di linguo internaciona\nl
en la maxim diversa landi e medii. Ni kun fido prizentas\nl
lu al publiko pro la alta kompetenteso di la ciencozi qui\nl
partoprenis en lua kompozo.
\end{VUpage}

% --- Leaf 10 (folio 6) ---------------------------------------------
%  Note rule: y=1527, 213 px, x0=196 for a margin at 195-199.
%  The detector did take the first line of BODY (y=247) and not the
%  folio (y=171, which would have given 139.10 mm).
%  The ten lines of note came out of the profile WELDED into four
%  ranges; separated at the baselines, not from the profile.
% -------------------------------------------------------------------
\begin{VUpage}[10]{6}
\VUnotes{133.14mm}{%
(1) Kontante segun la maniero dil Esperantisti, ni povus dicar,\nl
ke Ido havas nur 10 e mem 9 reguli, ne min kompleta. (Videz\nl
« Lexique-manuel », p. V-VIII.) --- \textit{Noto adjuntita.}

(2) En la realeso la gramatiko \textit{elementala} di Esperanto kontenas\nl
adminime 64 reguli e plu kam 10 ecepti. (Videz la konto detalizita\nl
ye la fino dil broshuro « \textit{La science et l'Esperanto} ».)

(3) Hike ni abandonas la Franca texto, qua deskriptas l'aranjeso\nl
di \textit{Grammaire complète}, e ni substitucas ad olu to, quo propre\nl
koncernas la « \textit{Kompleta Gramatiko detaloza} » quan ni prizentas\nl
ad omna amiki di helpolinguo.%
}

La kritiki, di qui ca gramatiko esos kredeble l'objekto,\nl
ne falios opozar ad olu la mikrega gramatiko primitiva di\nl
Esperanto, kun lua dek e sis reguli tante laudata (1). Ma,\nl
sen diskutar pri la numerizo di ta reguli, o plu \textit{vere para}\cc
\textit{grafi}, di qui un sola kontenas la konjugo (minus la formi\nl
kompozita per qui on mustis kompletigar lu pose), on darfas\nl
asertar sen ul paradoxo ke, se Esperanto havabus origine\nl
plu multa reguli, ol nun posedus quanto min granda de\nl
oli (2). Nam pro ta gramatiko vere tro kurta e nesuficanta,\nl
fakte a l'uzado Dro Zamenhof komisis la sorgo fixigar la\nl
formi di la linguo, ed ica uzado genitis multega partikularaji\nl
ed anomalaji, quin la gramatikisti di Esperanto kolektas e\nl
pie mencionas en lia verki, quale se traktesus linguo natu\cc
rala. E singla de ta partikularaji, singla de ta anomalaji\nl
efektigas un « regulo » quan la novici esas obligata lernar.\nl
Do esas tre vera dicar, ke se Esperanto havabus origine sat\nl
multa reguli generala e preciza, ol ne havus nun tro multa\nl
partikulara reguli, fondita sur « uzi » disparata e kontre\cc
logika. Pro to on obligesis kompozar fine grosa gramatiki,\nl
de qui un prizentas plu kam 400 paragrafi. A ca verki ya\nl
on devas equitatoze komparar nia \textit{Grammaire complète}. Lore\nl
on ne povos neagnoskar la supera simpleso e regulozeso di\nl
la linguo quan ni prizentas (3).

Certe ni povabus kondensar e multe plukurtigar ica verko,\nl
nome ne enduktar en lu kozi advere utila ma ne necesa, e\nl
qui ne esas strikte gramatikala. Ma lore ni agabus kontre\nl
nia intenci e l'explicita deziro di multa samideani. Cetere\nl
on dicernos facile la reguli lernenda del kozi nur lektebla\nl
o lektinda.

Kompreneble, en la gramatiko sancionita dal permananta
\parplein
\end{VUpage}

% --- Leaf 11 (folio 7) ---------------------------------------------
%  Note rule: y=1726, 212 px, x0=124 for a body at x0=121-126.
%  No centred rule, no ornament: nothing between 224 and 277, between
%  1686 and 1726, nor between 1916 and 2122.
% -------------------------------------------------------------------
\begin{VUpage}[11]{7}
\VUnotes{147.07mm}{%
(1) Ica vorti di nia unesma averto divenas partikulare vera pri\nl
la \textit{Kompleta Gramatiko detaloza,} en qua la multa pruntaji ek \textit{Pro}\cc
\textit{greso} e citaji di Sro \textsc{Couturat} posibligos ta profunda studiado.\nl
Danke oli, ni darfas dicar : mortinta, il parolas ankore, la granda\nl
kalumniato.%
}

\VUcontinue Komisitaro ni permisis a ni chanjar od adjuntar, pri la\nl
reguli, nur la kelka punti chanjita od adjuntita dal Ido-\nl
Akademio, legitima e yurizita sucedinto di olta. Ma la formo\nl
ed aranjeso necese diferas de olti, quin havis la verko pri\cc
mitiva, haste facita en tri monati e nur destinita al membri\nl
dil Komitato dil \textit{Delegitaro,} por utiligo eventuala. Certe la\nl
permananta Komisitaro ja emendabis ed igabis lu plu bona,\nl
ma vizante la linguo ipsa e ne olua docado specale. Or la\nl
publiko tre diferas de komitato di ciencisti, por qua multa\nl
kozi bezonis nek detali nek justifiko. Cetere la tempo mankis\nl
al Komisitaro por rilaborar detale la verko, nam adversi\nl
ed amiki ja esis incitanta a la propagado, iti kun intenco\nl
maligna, ici en sua fervoro difuzor lo plu bona. Rezulte, la\nl
primitiva gramatiko bezonis ankore plu kompleta riaranjo\nl
ed expliki plu longa. Ni esforcis en ica satisfacar ta duopla\nl
bezono e preferis donar tro multe kam ne sufice.

On permisez ad ni atraktar la atenco a kozo quan multi\nl
ne vidas sat bone. Pri la helpolinguo existas facileso trom\cc
panta, qua vizas nur o precipue \textit{lerno rapida} per nesuficanta\nl
reguli e nesuficanta radiki o vorti. Nu, ica facileso \textit{di lerno}\nl
ne povas konkordar e koincidar kun la facileso \textit{di apliko,} nek\nl
kun l'exakta e klara expresado dil pensi. Or on lernas la\nl
L. I. dum kelka hori, kelka jorni, ed on aplikas lu dum yari.\nl
Altraparte, quon valoras facileso e simpleso qui reale ne\nl
atingas la skopo? Mezurar la supereso, pri la helpolinguo,\nl
segun la minmulteso di la reguli e dil vorti, forsan esas\nl
habila trompilo por la maso, ma eroro deplorinda por omni,\nl
eroro qua produktas en la aplikado detrimenti maxim grava.\nl
Ido evitis fortunoze ta rifo mortigiva, kontre qua mult altri\nl
ja venis ed ankore venos ruptar su.

Nun ni finos, durante la traduko quan ni interruptis :

« Cetere, ica \textit{Kompleta Gramatiko} destinesas precipue al\nl
personi qui volas profundigar la studiado di la linguo (1).\nl
Ol esforcas solvar omna desfacilaji quin on povas renkontrar\parplein
\end{VUpage}

% --- Leaf 12 (folio 8): the end of the AVERTO ----------------------
%  NO BLOCK OF NOTES: « rule not found » is the right answer here.
%  The page's only stroke starts at x=648 and measures 197 px: it is
%  the END ORNAMENT, centred (middle 746 for a measure centred at
%  755), laid at 149.01 mm.
%  TWO SIGNATURES flush right, but drawn back by 41 and 43 px (3.47
%  and 3.64 mm): the same em on both -- hence \VUauFer.
%  « L. B. » HAD ESCAPED THE DETECTOR: 769 px of ink against 6882 for
%  a full line, that is 0.112 for a threshold set at 0.12. The page
%  carries 31 rows of ink, not 30.
%  « Sro » is in SMALL CAPITALS on this page, not superior as
%  everywhere else in the volume (witness: leaf 123).
%  « Chandelliez » is printed with two L's: kept.
% -------------------------------------------------------------------
\begin{VUpage}[12]{8}
\VUornement[16.68mm]{149.01mm}

\VUcontinue en l'uzado di ca linguo, e qui venas ne de ol ipsa, ma del\nl
neregulozaji e del idiotismi di nia lingui. Ta desfacilaji ren\cc
kontresas en omna traduko, e mem duopligita, kande on\nl
traktas naturala lingui, nam generale on esas obligata\nl
expresar l'idiotismi di un linguo per l'idiotismi di altra. Ma\nl
precize ta desfacilaji facas ek la L. I. bonega intelektal\nl
exerco, analoga (egardite omna proporciono) a la studio di\nl
la Latina o di la stranjera lingui :

Nam li kustumigas la spirito analizar la penso e liberigar\nl
ol del vesto plu o min trompiva, per qua nia lingui vesti\cc
zachas olu. Ma, vice quik impozar a lu nova travestio, la L. I.\nl
tendencas furnisar a lu l'expresuro maxim direta, maxim\nl
logikala e maxim klara. E, malgre l'opozata dici, logiko\nl
ya vere igas la L. I. supera a nia lingui « naturala ». Logiko\nl
ya sekurigas ad olu la palmo en lua aplikeso a la cienci;\nl
logiko donas ad olu l'astonanta facileso, kun qua omna\nl
populi ed omna klasi sociala povas komprenar olu, aquirar\nl
olu ed uzar olu.

\VUblanc{1.95mm}
\VUauFer{3.47mm}{\textsc{L. de Beaufront.}}
\VUblanc{3.13mm}

Me falius devo di yusteso e gratitudo, se me ne dankus\nl
hike \textsc{Sro Camille Chandelliez}, ex-docero dil grupo Espe\cc
rantista en \textit{Troyes}, Francia. Konvertita ad Ido per plur\cc
monata lekto e studio di \textit{Progreso}, il ofris a me sua serchi e\nl
noti, sua remarki e kritiki sagaca. Il helpis me tre efike\nl
klarigar ula punti, nome l'adjektivi-pronomi nedefinita. Ad\nl
ilu komplete debesas l'apendico pri la puntizado, la letri e\nl
adresi, segun \textit{Progreso}. Il esis mea unesma lektero e kon\cc
stanta kritikero. Ica gramatiko debas multo a lua helpado.\nl
Ton me joyas dicar publike e gratitudoze.

\VUblanc{1.95mm}
\VUauFer{3.64mm}{\textsc{L. B.}}
\end{VUpage}

% ===================================================================
%  LEAF 13 (folio 9) -- HALF TITLE OF THE FIRST PART.
%  Its exact precedent is leaf 119, the half title of the second part:
%  the same sizes (14.17 pt at the large line in both cases), the same
%  intervals to within 0.25 mm, less the subtitle line.
%
%  The page has neither folio nor line of body. Calibration used:
%  leaf 7's -- (top_of_capitals - edge_of_paper) x 25,4/300 + C, with
%  C = 5.84 mm. The surveyor RECOMPUTED C over 34 pages with a visible
%  edge: median 5.82 mm, standard deviation 1.15 -- the project's
%  constant is confirmed.
%  AN INDEPENDENT CROSS-CHECK: leaf 12 comes from the SAME raw scan
%  (the left half of the same shot) and does carry a folio and text.
%  Going by its first line of body (y=275), the volume's ordinary rule
%  gives 72.31 / 81.79 / 90.76 mm. The two calibrations agree to
%  within 0.21 mm. It is the first that are kept.
%
%  THE TWO LINES ARE ROMAN, not bold: median stem 5 px for a capital
%  of 29 (0.172) and 7 px for 40 (0.175), where the set semibold gives
%  10. There is no roman on these lines to serve as a witness -- they
%  are alone on the page -- so that the comparison is with the roman
%  SET at the same size. The method was checked by applying it again
%  to leaf 119, where it gives back the figures already recorded.
%
%  The letterspacing is CORRECTED ON THE SET TEXT, not taken from the
%  bare rule of apparatus.py, which gave 111 for the large line
%  (2.5 mm too short). The values laid give 216 and 813 px, that is
%  the facsimile's measurements to the pixel.
%
%  A RESERVATION BEARING ON ANOTHER PAGE: the surveyor holds, with
%  measurements to support it, that LEAF 119 is laid 7.7 mm too high,
%  its ordinates being the raw y x 25,4/300 of apparatus.py, which
%  assumes that the first row of the crop is the edge of the paper --
%  false on that shot. See docs/journal.md, the section of open
%  checks: the reservation is recorded, it is not yet verified, and
%  leaf 119 has not been touched.
% ===================================================================
\begin{VUpage}[13]{}
\VUcentreA{72.09mm}{10.27pt}{182}{I\textsuperscript{a} PARTO}
\VUfiletA{81.57mm}{8.30mm}
\VUcentreA{90.55mm}{14.17pt}{134}{MORFOLOGIO E SINTAXO}
\end{VUpage}

% --- Leaf 14 (folio 10): blank -------------------------------------
%  225 px of ink at threshold 140: the page is blank.
\begin{VUpage}[14]{}
\end{VUpage}
